\IfFileExists{stacks-project.cls}{%
\documentclass{stacks-project}
}{%
\documentclass{amsart}
}

% For dealing with references we use the comment environment
\usepackage{verbatim}
\newenvironment{reference}{\comment}{\endcomment}
%\newenvironment{reference}{}{}
\newenvironment{slogan}{\comment}{\endcomment}
\newenvironment{history}{\comment}{\endcomment}

% For commutative diagrams we use Xy-pic
\usepackage[all]{xy}

% We use 2cell for 2-commutative diagrams.
\xyoption{2cell}
\UseAllTwocells

% We use multicol for the list of chapters between chapters
\usepackage{multicol}

% This is generally recommended for better output
\usepackage{lmodern}
\usepackage[T1]{fontenc}

% For cross-file-references
\usepackage{xr-hyper}

% Package for hypertext links:
\usepackage{hyperref}

% For any local file, say "hello.tex" you want to link to please
% use \externaldocument[hello-]{hello}
\externaldocument[introduction-]{introduction}
\externaldocument[conventions-]{conventions}
\externaldocument[sets-]{sets}
\externaldocument[categories-]{categories}
\externaldocument[topology-]{topology}
\externaldocument[sheaves-]{sheaves}
\externaldocument[sites-]{sites}
\externaldocument[stacks-]{stacks}
\externaldocument[fields-]{fields}
\externaldocument[algebra-]{algebra}
\externaldocument[brauer-]{brauer}
\externaldocument[homology-]{homology}
\externaldocument[derived-]{derived}
\externaldocument[simplicial-]{simplicial}
\externaldocument[more-algebra-]{more-algebra}
\externaldocument[smoothing-]{smoothing}
\externaldocument[modules-]{modules}
\externaldocument[sites-modules-]{sites-modules}
\externaldocument[injectives-]{injectives}
\externaldocument[cohomology-]{cohomology}
\externaldocument[sites-cohomology-]{sites-cohomology}
\externaldocument[dga-]{dga}
\externaldocument[dpa-]{dpa}
\externaldocument[sdga-]{sdga}
\externaldocument[hypercovering-]{hypercovering}
\externaldocument[schemes-]{schemes}
\externaldocument[constructions-]{constructions}
\externaldocument[properties-]{properties}
\externaldocument[morphisms-]{morphisms}
\externaldocument[coherent-]{coherent}
\externaldocument[divisors-]{divisors}
\externaldocument[limits-]{limits}
\externaldocument[varieties-]{varieties}
\externaldocument[topologies-]{topologies}
\externaldocument[descent-]{descent}
\externaldocument[perfect-]{perfect}
\externaldocument[more-morphisms-]{more-morphisms}
\externaldocument[flat-]{flat}
\externaldocument[groupoids-]{groupoids}
\externaldocument[more-groupoids-]{more-groupoids}
\externaldocument[etale-]{etale}
\externaldocument[chow-]{chow}
\externaldocument[intersection-]{intersection}
\externaldocument[pic-]{pic}
\externaldocument[weil-]{weil}
\externaldocument[adequate-]{adequate}
\externaldocument[dualizing-]{dualizing}
\externaldocument[duality-]{duality}
\externaldocument[discriminant-]{discriminant}
\externaldocument[derham-]{derham}
\externaldocument[local-cohomology-]{local-cohomology}
\externaldocument[algebraization-]{algebraization}
\externaldocument[curves-]{curves}
\externaldocument[resolve-]{resolve}
\externaldocument[models-]{models}
\externaldocument[functors-]{functors}
\externaldocument[equiv-]{equiv}
\externaldocument[pione-]{pione}
\externaldocument[etale-cohomology-]{etale-cohomology}
\externaldocument[proetale-]{proetale}
\externaldocument[relative-cycles-]{relative-cycles}
\externaldocument[more-etale-]{more-etale}
\externaldocument[trace-]{trace}
\externaldocument[crystalline-]{crystalline}
\externaldocument[spaces-]{spaces}
\externaldocument[spaces-properties-]{spaces-properties}
\externaldocument[spaces-morphisms-]{spaces-morphisms}
\externaldocument[decent-spaces-]{decent-spaces}
\externaldocument[spaces-cohomology-]{spaces-cohomology}
\externaldocument[spaces-limits-]{spaces-limits}
\externaldocument[spaces-divisors-]{spaces-divisors}
\externaldocument[spaces-over-fields-]{spaces-over-fields}
\externaldocument[spaces-topologies-]{spaces-topologies}
\externaldocument[spaces-descent-]{spaces-descent}
\externaldocument[spaces-perfect-]{spaces-perfect}
\externaldocument[spaces-more-morphisms-]{spaces-more-morphisms}
\externaldocument[spaces-flat-]{spaces-flat}
\externaldocument[spaces-groupoids-]{spaces-groupoids}
\externaldocument[spaces-more-groupoids-]{spaces-more-groupoids}
\externaldocument[bootstrap-]{bootstrap}
\externaldocument[spaces-pushouts-]{spaces-pushouts}
\externaldocument[spaces-chow-]{spaces-chow}
\externaldocument[groupoids-quotients-]{groupoids-quotients}
\externaldocument[spaces-more-cohomology-]{spaces-more-cohomology}
\externaldocument[spaces-simplicial-]{spaces-simplicial}
\externaldocument[spaces-duality-]{spaces-duality}
\externaldocument[formal-spaces-]{formal-spaces}
\externaldocument[restricted-]{restricted}
\externaldocument[spaces-resolve-]{spaces-resolve}
\externaldocument[formal-defos-]{formal-defos}
\externaldocument[defos-]{defos}
\externaldocument[cotangent-]{cotangent}
\externaldocument[examples-defos-]{examples-defos}
\externaldocument[algebraic-]{algebraic}
\externaldocument[examples-stacks-]{examples-stacks}
\externaldocument[stacks-sheaves-]{stacks-sheaves}
\externaldocument[criteria-]{criteria}
\externaldocument[artin-]{artin}
\externaldocument[quot-]{quot}
\externaldocument[stacks-properties-]{stacks-properties}
\externaldocument[stacks-morphisms-]{stacks-morphisms}
\externaldocument[stacks-limits-]{stacks-limits}
\externaldocument[stacks-cohomology-]{stacks-cohomology}
\externaldocument[stacks-perfect-]{stacks-perfect}
\externaldocument[stacks-introduction-]{stacks-introduction}
\externaldocument[stacks-more-morphisms-]{stacks-more-morphisms}
\externaldocument[stacks-geometry-]{stacks-geometry}
\externaldocument[moduli-]{moduli}
\externaldocument[moduli-curves-]{moduli-curves}
\externaldocument[examples-]{examples}
\externaldocument[exercises-]{exercises}
\externaldocument[guide-]{guide}
\externaldocument[desirables-]{desirables}
\externaldocument[coding-]{coding}
\externaldocument[obsolete-]{obsolete}
\externaldocument[fdl-]{fdl}
\externaldocument[index-]{index}

% Theorem environments.
%
\theoremstyle{plain}
\newtheorem{theorem}[subsection]{Theorem}
\newtheorem{proposition}[subsection]{Proposition}
\newtheorem{lemma}[subsection]{Lemma}

\theoremstyle{definition}
\newtheorem{definition}[subsection]{Definition}
\newtheorem{example}[subsection]{Example}
\newtheorem{exercise}[subsection]{Exercise}
\newtheorem{situation}[subsection]{Situation}

\theoremstyle{remark}
\newtheorem{remark}[subsection]{Remark}
\newtheorem{remarks}[subsection]{Remarks}

\numberwithin{equation}{subsection}

% Macros
%
\def\lim{\mathop{\mathrm{lim}}\nolimits}
\def\colim{\mathop{\mathrm{colim}}\nolimits}
\def\Spec{\mathop{\mathrm{Spec}}}
\def\Hom{\mathop{\mathrm{Hom}}\nolimits}
\def\Ext{\mathop{\mathrm{Ext}}\nolimits}
\def\SheafHom{\mathop{\mathcal{H}\!\mathit{om}}\nolimits}
\def\SheafExt{\mathop{\mathcal{E}\!\mathit{xt}}\nolimits}
\def\Sch{\mathit{Sch}}
\def\Mor{\mathop{\mathrm{Mor}}\nolimits}
\def\Ob{\mathop{\mathrm{Ob}}\nolimits}
\def\Sh{\mathop{\mathit{Sh}}\nolimits}
\def\NL{\mathop{N\!L}\nolimits}
\def\CH{\mathop{\mathrm{CH}}\nolimits}
\def\proetale{{pro\text{-}\acute{e}tale}}
\def\etale{{\acute{e}tale}}
\def\QCoh{\mathit{QCoh}}
\def\Ker{\mathop{\mathrm{Ker}}}
\def\Im{\mathop{\mathrm{Im}}}
\def\Coker{\mathop{\mathrm{Coker}}}
\def\Coim{\mathop{\mathrm{Coim}}}

% Boxtimes
%
\DeclareMathSymbol{\boxtimes}{\mathbin}{AMSa}{"02}

%
% Macros for moduli stacks/spaces
%
\def\QCohstack{\mathcal{QC}\!\mathit{oh}}
\def\Cohstack{\mathcal{C}\!\mathit{oh}}
\def\Spacesstack{\mathcal{S}\!\mathit{paces}}
\def\Quotfunctor{\mathrm{Quot}}
\def\Hilbfunctor{\mathrm{Hilb}}
\def\Curvesstack{\mathcal{C}\!\mathit{urves}}
\def\Polarizedstack{\mathcal{P}\!\mathit{olarized}}
\def\Complexesstack{\mathcal{C}\!\mathit{omplexes}}
% \Pic is the operator that assigns to X its picard group, usage \Pic(X)
% \Picardstack_{X/B} denotes the Picard stack of X over B
% \Picardfunctor_{X/B} denotes the Picard functor of X over B
\def\Pic{\mathop{\mathrm{Pic}}\nolimits}
\def\Picardstack{\mathcal{P}\!\mathit{ic}}
\def\Picardfunctor{\mathrm{Pic}}
\def\Deformationcategory{\mathcal{D}\!\mathit{ef}}
\usepackage{cleveref}
\crefname{originalreference}{original reference}{original references}
\Crefname{originalreference}{Original reference}{Original references}
\makeatletter
\newlabel{algebra-lemma-Noetherian-irreducible-components}{{00FR}{source}{}{sourcepointer-1}{}}
\newlabel{algebra-lemma-Noetherian-irreducible-components@cref}{{[originalreference][0][]00FR}{[1][0][]source}{}{}{}}
\makeatother
\crefname{originalreference}{original reference}{original references}
\Crefname{originalreference}{Original reference}{Original references}
\makeatletter
\newlabel{algebra-lemma-total-ring-fractions-no-embedded-points}{{02LX}{source}{}{sourcepointer-2}{}}
\newlabel{algebra-lemma-total-ring-fractions-no-embedded-points@cref}{{[originalreference][0][]02LX}{[1][0][]source}{}{}{}}
\makeatother
\crefname{originalreference}{original reference}{original references}
\Crefname{originalreference}{Original reference}{Original references}
\makeatletter
\newlabel{algebra-lemma-when-colimit}{{07C3}{source}{}{sourcepointer-3}{}}
\newlabel{algebra-lemma-when-colimit@cref}{{[originalreference][0][]07C3}{[1][0][]source}{}{}{}}
\makeatother
\crefname{originallemma}{original lemma}{original lemmas}
\Crefname{originallemma}{Original lemma}{Original lemmas}
\makeatletter
\newlabel{lemma-delocalize-base}{{07F4}{source}{}{sourcepointer-4}{}}
\newlabel{lemma-delocalize-base@cref}{{[originallemma][0][]07F4}{[1][0][]source}{}{}{}}
\makeatother
\crefname{originallemma}{original lemma}{original lemmas}
\Crefname{originallemma}{Original lemma}{Original lemmas}
\makeatletter
\newlabel{lemma-desingularize-strictly-standard}{{07CT}{source}{}{sourcepointer-5}{}}
\newlabel{lemma-desingularize-strictly-standard@cref}{{[originallemma][0][]07CT}{[1][0][]source}{}{}{}}
\makeatother
\crefname{originallemma}{original lemma}{original lemmas}
\Crefname{originallemma}{Original lemma}{Original lemmas}
\makeatletter
\newlabel{lemma-lifting}{{07CP}{source}{}{sourcepointer-6}{}}
\newlabel{lemma-lifting@cref}{{[originallemma][0][]07CP}{[1][0][]source}{}{}{}}
\makeatother
\crefname{originallemma}{original lemma}{original lemmas}
\Crefname{originallemma}{Original lemma}{Original lemmas}
\makeatletter
\newlabel{lemma-product}{{07F3}{source}{}{sourcepointer-7}{}}
\newlabel{lemma-product@cref}{{[originallemma][0][]07F3}{[1][0][]source}{}{}{}}
\makeatother
\crefname{originalreference}{original reference}{original references}
\Crefname{originalreference}{Original reference}{Original references}
\makeatletter
\newlabel{more-algebra-lemma-regular-base-change}{{07C1}{source}{}{sourcepointer-8}{}}
\newlabel{more-algebra-lemma-regular-base-change@cref}{{[originalreference][0][]07C1}{[1][0][]source}{}{}{}}
\makeatother
\crefname{originalproposition}{original proposition}{original propositions}
\Crefname{originalproposition}{Original proposition}{Original propositions}
\makeatletter
\newlabel{proposition-lift}{{07CM}{source}{}{sourcepointer-9}{}}
\newlabel{proposition-lift@cref}{{[originalproposition][0][]07CM}{[1][0][]source}{}{}{}}
\makeatother
\crefname{originalreference}{original reference}{original references}
\Crefname{originalreference}{Original reference}{Original references}
\makeatletter
\newlabel{situation-global}{{07F2}{source}{}{sourcepointer-10}{}}
\newlabel{situation-global@cref}{{[originalreference][0][]07F2}{[1][0][]source}{}{}{}}
\makeatother
\title{The approximation property reduces to algebras geometrically regular over a field}
\author{The Stacks Project Authors}
\date{stacks-project:smoothing.tex:local-snapshot-20261012; GNU FDL 1.2 or later}
\begin{document}
\maketitle
\noindent\textit{Editorial note.} Input lemma-product is same-chapter NOT carded (dropped in judgment: Omitted with Hint) and is declared external; input lemma-lifting is published (tag 07CP).\par\medskip
\noindent\textit{Allowed standard prerequisites (editorial).}\begin{itemize}
\item Regular maps are stable under quotient base change. More on Algebra, lemma regular-base-change (cross-chapter, named). Keeps R/I -> Lambda/ILambda regular along the Noetherian induction.
\item Ind-smoothness is stable under infinitesimal deformations. proposition-lift (carded L26807). Shows the maximal counterexample ring R is reduced.
\item Products of PT situations. lemma-product (same chapter, NOT carded, declared external). Transfers PT from the fields K\_1 x ... x K\_n to S\textasciicircum{}\{-1\}R -> S\textasciicircum{}\{-1\}Lambda.
\item Spreading a smooth factorization out of a localization. lemma-delocalize-base (carded L26810). Produces the factorization with an elementary standard denominator pi in S.
\item Lifting along annihilator-controlled elements. lemma-lifting (same chapter, PUBLISHED 07CP). Lifts the mod-pi\textasciicircum{}8 factorization to a smooth D over R with a mod-pi\textasciicircum{}4 factorization.
\item Desingularization with compatible mod-pi\textasciicircum{}4 maps. lemma-desingularize-strictly-standard (carded L26808). Final step: converts the mod-pi\textasciicircum{}4 compatibility into an honest smooth factorization.
\end{itemize}\medskip
\ifdefined\hypertarget\hypertarget{sourcepointer-1}{}\fi
\noindent\textit{Source reference.} Original reference 00FR: focused-seeds/geometry/S9-smo-reduce-to-field/ref-tags.txt, source lines 5--5. Navigation pointer only: the counted proof cites this statement by \textbackslash{}ref as a standing input; the deduction the counted proof has to make is carried out inside the counted span, and no part of that deduction is delegated to this pointer.\par\medskip
\ifdefined\hypertarget\hypertarget{sourcepointer-2}{}\fi
\noindent\textit{Source reference.} Original reference 02LX: focused-seeds/geometry/S9-smo-reduce-to-field/ref-tags.txt, source lines 6--6. Navigation pointer only: the counted proof cites this statement by \textbackslash{}ref as a standing input; the deduction the counted proof has to make is carried out inside the counted span, and no part of that deduction is delegated to this pointer.\par\medskip
\ifdefined\hypertarget\hypertarget{sourcepointer-3}{}\fi
\noindent\textit{Source reference.} Original reference 07C3: focused-seeds/geometry/S9-smo-reduce-to-field/ref-tags.txt, source lines 7--7. Navigation pointer only: the counted proof cites this statement by \textbackslash{}ref as a standing input; the deduction the counted proof has to make is carried out inside the counted span, and no part of that deduction is delegated to this pointer.\par\medskip
\ifdefined\hypertarget\hypertarget{sourcepointer-4}{}\fi
\noindent\textit{Source reference.} Original lemma 07F4: sources/stacks-smoothing/pinned-original-source.tex, source lines 2154--2154. Navigation pointer only: the counted proof cites this statement by \textbackslash{}ref as a standing input; the deduction the counted proof has to make is carried out inside the counted span, and no part of that deduction is delegated to this pointer.\par\medskip
\ifdefined\hypertarget\hypertarget{sourcepointer-5}{}\fi
\noindent\textit{Source reference.} Original lemma 07CT: sources/stacks-smoothing/pinned-original-source.tex, source lines 2029--2029. Navigation pointer only: the counted proof cites this statement by \textbackslash{}ref as a standing input; the deduction the counted proof has to make is carried out inside the counted span, and no part of that deduction is delegated to this pointer.\par\medskip
\ifdefined\hypertarget\hypertarget{sourcepointer-6}{}\fi
\noindent\textit{Source reference.} Original lemma 07CP: sources/stacks-smoothing/pinned-original-source.tex, source lines 1638--1638. Navigation pointer only: the counted proof cites this statement by \textbackslash{}ref as a standing input; the deduction the counted proof has to make is carried out inside the counted span, and no part of that deduction is delegated to this pointer.\par\medskip
\ifdefined\hypertarget\hypertarget{sourcepointer-7}{}\fi
\noindent\textit{Source reference.} Original lemma 07F3: sources/stacks-smoothing/pinned-original-source.tex, source lines 2143--2143. Navigation pointer only: the counted proof cites this statement by \textbackslash{}ref as a standing input; the deduction the counted proof has to make is carried out inside the counted span, and no part of that deduction is delegated to this pointer.\par\medskip
\ifdefined\hypertarget\hypertarget{sourcepointer-8}{}\fi
\noindent\textit{Source reference.} Original reference 07C1: focused-seeds/geometry/S9-smo-reduce-to-field/ref-tags.txt, source lines 8--8. Navigation pointer only: the counted proof cites this statement by \textbackslash{}ref as a standing input; the deduction the counted proof has to make is carried out inside the counted span, and no part of that deduction is delegated to this pointer.\par\medskip
\ifdefined\hypertarget\hypertarget{sourcepointer-9}{}\fi
\noindent\textit{Source reference.} Original proposition 07CM: sources/stacks-smoothing/pinned-original-source.tex, source lines 1591--1591. Navigation pointer only: the counted proof cites this statement by \textbackslash{}ref as a standing input; the deduction the counted proof has to make is carried out inside the counted span, and no part of that deduction is delegated to this pointer.\par\medskip
\ifdefined\hypertarget\hypertarget{sourcepointer-10}{}\fi
\noindent\textit{Source reference.} Original reference 07F2: sources/stacks-smoothing/pinned-original-source.tex, source lines 2133--2133. Navigation pointer only: the counted proof cites this statement by \textbackslash{}ref as a standing input; the deduction the counted proof has to make is carried out inside the counted span, and no part of that deduction is delegated to this pointer.\par\medskip
\begin{lemma}
\label{lemma-reduce-to-field}
\begin{slogan}
Proving Popescu approximation reduces to algebras over a field
\end{slogan}
If for every Situation \ref{situation-global} where $R$
is a field PT holds, then PT holds in general.
\end{lemma}

\begin{proof}
Assume PT holds for any Situation \ref{situation-global} where $R$ is a field.
Let $R \to \Lambda$ be as in Situation \ref{situation-global} arbitrary.
Note that $R/I \to \Lambda/I\Lambda$ is another regular ring map
of Noetherian rings, see
More on Algebra, Lemma \ref{more-algebra-lemma-regular-base-change}.
Consider the set of ideals
$$
\mathcal{I} = \{I \subset R \mid R/I \to \Lambda/I\Lambda
\text{ does not have PT}\}
$$
We have to show that $\mathcal{I}$ is empty. If this set is nonempty,
then it contains a maximal element because $R$ is Noetherian.
Replacing $R$ by $R/I$ and $\Lambda$ by $\Lambda/I$ we obtain a
situation where PT holds for $R/I \to \Lambda/I\Lambda$ for any
nonzero ideal of $R$. In particular, we see by applying
Proposition \ref{proposition-lift}
that $R$ is a reduced ring.

\medskip\noindent
Let $A \to \Lambda$ be an $R$-algebra homomorphism with $A$ of
finite presentation. We have to find a factorization $A \to B \to \Lambda$
with $B$ smooth over $R$, see Algebra, Lemma \ref{algebra-lemma-when-colimit}.

\medskip\noindent
Let $S \subset R$ be the set of nonzerodivisors and
consider the total ring of fractions $Q = S^{-1}R$ of $R$. We know that
$Q = K_1 \times \ldots \times K_n$ is a product of fields, see
Algebra, Lemmas \ref{algebra-lemma-total-ring-fractions-no-embedded-points} and
\ref{algebra-lemma-Noetherian-irreducible-components}.
By Lemma \ref{lemma-product} and our assumption
PT holds for the ring map $S^{-1}R \to S^{-1}\Lambda$.
Hence we can find a factorization $S^{-1}A \to B' \to S^{-1}\Lambda$
with $B'$ smooth over $S^{-1}R$.

\medskip\noindent
We apply Lemma \ref{lemma-delocalize-base}
and find a factorization $A \to B \to \Lambda$ such that
some $\pi \in S$ is elementary standard in $B$ over $R$.
After replacing $A$ by $B$ we may assume that $\pi$ is
elementary standard, hence strictly standard in $A$. We know that
$R/\pi^8R \to \Lambda/\pi^8\Lambda$ satisfies PT.
Hence we can find a factorization
$R/\pi^8 R \to A/\pi^8A \to \bar C \to \Lambda/\pi^8\Lambda$
with $R/\pi^8 R \to \bar C$ smooth. By
Lemma \ref{lemma-lifting}
we can find an $R$-algebra map $D \to \Lambda$ with $D$ smooth over $R$
and a factorization
$R/\pi^4 R \to A/\pi^4A \to D/\pi^4D \to \Lambda/\pi^4\Lambda$.
By Lemma \ref{lemma-desingularize-strictly-standard}
we can find $A \to B \to \Lambda$ with $B$ smooth over $R$
which finishes the proof.
\end{proof}

\end{document}
